\bookchapter{บทที่ 50: เริ่มต้น Kotlin อย่างรวดเร็ว — ภาษาหลักของการพัฒนา Android}

\chapterintro{บทนี้เราจะเริ่มเรียนภาษากันจริงจัง ไม่ต้องเครียด — Kotlin ขึ้นชื่อว่า “เป็นมิตรกับมือใหม่” แถมคุณไม่ต้องท่องหนังสือไวยากรณ์: แค่เข้าใจไวยากรณ์ 20\% แรก ก็อ่านรู้เรื่องและเขียนโค้ดประจำวันได้ถึง 80\% ส่วนที่เหลือค่อยค้นระหว่างทำโปรเจกต์ไป จำหลักการเรียนของเราไว้: เข้าใจ > ท่องจำ ใช้เป็น > ชำนาญ}

\booksection{ตัวแปรกับชนิดข้อมูล: ตั้งชื่อให้ข้อมูล}

โปรแกรมก็คือ “ข้อมูล + ตรรกะ” ตัวแปร ก็คือการตั้งชื่อให้ข้อมูลแล้วเก็บไว้ Kotlin มีสองวิธีประกาศ:

\begin{lstlisting}[language=Kotlin]
val name = "ซานผวน"        // val: ตัวแปรอ่านอย่างเดียว (แนะนำให้ใช้เป็นค่าเริ่มต้น)
var age = 18             // var: ตัวแปรแก้ได้ (ภายหลังเปลี่ยนได้)

age = 19                 // <i class="fas fa-check-circle mr-1"></i> var แก้ได้
// name = "ลี่ซือ"         // <i class="fas fa-times-circle mr-1"></i> val แก้ไม่ได้ คอมไพล์แล้ว error ทันที

val score: Int = 100     // จะเขียนไทป์ชัด ๆ ก็ได้
val pi = 3.14            // ไม่เขียนก็เดาเองได้
\end{lstlisting}

ชนิดที่ใช้บ่อยมีแค่ไม่กี่ตัว จำจากตารางเดียวจบ:

\begin{center}
\begin{tabularx}{\linewidth}{|X|X|X|}
\hline
\textbf{ชนิด} & \textbf{เก็บอะไร} & \textbf{ตัวอย่าง} \tabularnewline
\hline
Int & จำนวนเต็ม & 42 \tabularnewline
Double & ทศนิยม & 3.14 \tabularnewline
String & ข้อความ & "สวัสดี" \tabularnewline
Boolean & จริง/เท็จ & true / false \tabularnewline
List< T> & ชุดข้อมูล & listOf(1, 2, 3) \tabularnewline
\hline
\end{tabularx}
\end{center}

\begin{center}\small\color{gray}ชนิดข้อมูลพื้นฐานที่ใช้บ่อยใน Kotlin\end{center}

\begin{tcbinfo}{เคล็ดลับ: หลัก val มาก่อน}
กฎเหล็กของชุมชน Kotlin: ใช้ val เป็นค่าเริ่มต้น จะใช้ var เมื่อจำเป็นจริง ๆ ถ้าที่ที่ใช้ val ได้แล้วไปใช้ var IDE จะเตือนด้วยตัวเหลือง นิสัยนี้ช่วยตัดบั๊กกลุ่มใหญ่ “ตัวแปรถูกแก้โดยไม่ตั้งใจ” ทิ้งได้ — โค้ดยิ่งมีจุดเปลี่ยนน้อย ยิ่งไล่เหตุผลง่าย

\end{tcbinfo}

\booksection{ฟังก์ชัน: รวมตรรกะเป็น “ปุ่ม” พร้อมกด}

ฟังก์ชันคือตรรกะชิ้นหนึ่งที่มีชื่อ เหมือนปุ่มบนรีโมต — กดหนึ่งที (เรียกใช้) ก็ทำชุดการกระทำต่อให้ทั้งหมด ใน Kotlin ประกาศด้วย fun:

\begin{lstlisting}[language=Kotlin]
fun greet(name: String): String {
    return "สวัสดี $name!"     // เครื่องหมาย $ สอดตัวแปรลงในข้อความ
}

println(greet("ซานผวน"))        // ผลลัพธ์: สวัสดี ซานผวน!

// ฟังก์ชันบรรทัดเดียวเขียนเป็น expression ได้
fun add(a: Int, b: Int) = a + b

println(add(3, 5))            // ผลลัพธ์: 8
\end{lstlisting}

จุดที่ต้องรู้มีดังนี้: พารามิเตอร์เขียนในวงเล็บ (ชื่อมาก่อน ชนิดมาหลัง คั่นด้วยทวิภาค); ชนิดค่าที่ส่งกลับเขียนหลังวงเล็บ; สตริงเทมเพลต \$ชื่อตัวแปร คือของเด่นสะดวกสบายของ Kotlin — ไม่ต้องเอา “สตริง + ตัวแปร” มาต่อกันแบบ Java อีกแล้ว

\booksection{เงื่อนไขกับลูป: ทำให้โปรแกรม “คิด” เป็น}

ความฉลาดของโปรแกรมมาจากสองเรื่อง: ตัดสินใจเป็น (if/when) กับทำซ้ำเป็น (for/while)

\begin{lstlisting}[language=Kotlin]
// การเช็คเงื่อนไข: if / else
val age = 17
if (age >= 18) {
    println("ผู้ใหญ่")
} else {
    println("ผู้เยาว์")
}

// หลายทางเลือก: when (ใช้ดีกว่า switch ร้อยเท่า)
val score = 85
val level = when {
    score >= 90 -> "ดีเยี่ยม"
    score >= 60 -> "ผ่าน"
    else -> "สู้ ๆ เข้าละ!"
}

// ทำซ้ำ: ลูป for
val fruits = listOf("แอปเปิล", "กล้วย", "ส้ม")
for (fruit in fruits) {
    println("ฉันชอบกิน$fruit")
}
\end{lstlisting}

when คือรูปแบบหลายทางเลือกที่ใช้บ่อยที่สุดใน Kotlin งานพัฒนา Android คุณจะได้ใช้มันทุกวัน (เช่น แสดงหน้าจอต่างกันตามสถานะ)

\booksection{Null Safety: “เครื่องรางกันตาย” ของ Kotlin}

บั๊กที่ดังที่สุดยุค Java ชื่อ NullPointerException — โปรแกรมเผลอไปเข้าถึงสิ่งที่ “ไม่มีอะไรเลย” แล้วดับทันที ตามสถิติมันครองอันดับหนึ่งบนกระดานแอป Android พังมาหลายปีติดต่อกัน Kotlin แก้ปัญหานี้ตั้งแต่ระดับภาษา: โดยปริยายห้ามค่า null

\begin{lstlisting}[language=Kotlin]
var name: String = "ซานผวน"   // ไทป์ปกติ: ห้ามเป็น null เด็ดขาด
// name = null                // <i class="fas fa-times-circle mr-1"></i> แก้ก่อน ไม่งั้นคอมไพเลอร์ตีกลับทันที

var nickname: String? = null // ใส่เครื่องหมายคำถาม: อนุญาตให้เป็น null
nickname = "ซานเอ๋อร์"

// ใช้ nullable type ต้องจัดการอย่างปลอดภัย:
println(nickname?.length)    // safe call: เป็น null ก็คืน null ไม่เด้ง
val len = nickname?.length ?: 0   // Elvis operator: เป็น null ก็ให้ค่าเริ่มต้น
println(len)
\end{lstlisting}

\begin{tcbwarn}{ข้อควรรู้: อย่าสับสนระหว่าง ? กับ ?}
String? คือ “อนุญาตให้เป็น null ได้” ส่วน a ?: b คือ “ถ้า a เป็น null ให้ใช้ b” (หน้าตาเหมือนทรงผมเอลวิส เลยเรียกกันว่า Elvis operator) ช่วงแรกที่เขียนโค้ด พอเห็นคอมไพเลอร์ฟ้อง “error เรื่อง null” แปดทับสิบคือคุณลืมจัดการชนิดที่เป็น null ได้ — นั่นแหละคือภาษากำลังคุ้มครองคุณอยู่

\end{tcbwarn}

\booksection{คลาสกับอ็อบเจกต์: แม่พิมพ์สำหรับ “หล่อชิ้นส่วน”}

คลาสคือ “แม่พิมพ์” ของอ็อบเจกต์ อ็อบเจกต์คือ “ชิ้นงานจริง” ที่หล่อจากคลาส ทุกอย่างใน Android — หน้าจอ ปุ่ม ข้อมูล — ล้วนเป็นอ็อบเจกต์ ใน Kotlin มีของวิเศษชื่อ data class (คลาสข้อมูล) ทำหน้าที่ “ใส่ข้อมูล” โดยเฉพาะ หนึ่งบรรทัดแทน Java ได้หลายสิบบรรทัด:

\begin{lstlisting}[language=Kotlin]
// คลาสปกติ: แม่พิมพ์
class Dog(val name: String) {
    fun bark() = println("$name ฮ่อง ฮ่อง!")
}

val dog = Dog("โอ้เจ้า")      // สร้างออบเจกต์
dog.bark()                 // โอ้เจ้า ฮ่อง ฮ่อง!

// data class: กล่องใส่ข้อมูล (มี equals/copy ให้พร้อม)
data class User(val id: Int, val name: String, val email: String)

val me = User(1, "ซานผวน", "zhangsan@example.com")
val twin = me.copy(name = "ซานซานเฟิง")   // คัดลอกแล้วแก้ค่าเดียว
\end{lstlisting}

ต่อไปถ้าคุณทำแอปบันทึกบัญชี หนึ่งรายการก็คือ data class หนึ่งตัว: data class Record(val amount: Double, val category: String, val date: String) เห็นไหมล่ะ คุณเขียนโครงสร้างโค้ดระดับโปรเจกต์จริงได้แล้ว

\booksection{lambda กับ high-order function: “น้ำตาลไวยากรณ์” ของ Kotlin}

ในโค้ด Android คุณจะเจอ “ฟังก์ชันไร้ชื่อ” แบบ \{ ... \} ทุกวัน มันชื่อ lambda (ฟังก์ชันนิรนาม) onClick = \{ count++ \} ของ Button ก็คือ lambda หนึ่งตัว — ห่อ “กดแล้วทำอะไร” เป็นแพ็กเดียว แล้วส่งต่อให้ปุ่มไป

\begin{lstlisting}[language=Kotlin]
// แบบปกติ: นิยามฟังก์ชันมีชื่อก่อนแล้วค่อยส่งเข้าไป
fun double(x: Int) = x * 2
val result = listOf(1, 2, 3).map(::double)

// สไตล์ lambda: นิยามเลยตรงนั้น กระชับไว
val result2 = listOf(1, 2, 3).map { it * 2 }   // it = ตัวปัจจุบัน

// กรอง + แปลงต่อกัน (เจอบ่อยมากในโปรเจกต์จริง)
val bills = listOf(12.5, 88.0, 5.0, 200.0)
val bigBills = bills.filter { it > 50 }.sorted()
// [88.0, 200.0]—"บิลที่เกิน 50 เรียงให้เรียบร้อย"
\end{lstlisting}

จำคอมโบ “คอลเลกชัน” ที่ใช้บ่อยไว้หน่อย: filter (กรอง), map (แปลง), sorted (เรียง), sum (รวม) งานจัดการลิสต์ของคุณในอนาคต 90\% คือการผสมสี่ท่านี้ ปรัชญาของ Kotlin คือ: ตรรกะที่อธิบายให้ชัดได้ในหนึ่งประโยค ห้ามเขียนกินห้าประโยค

\booksection{การ์ดสรุปไวยากรณ์ฉบับคนเริ่มจากศูนย์}

อย่าหวังท่องไวยากรณ์ทั้งหมดให้จบในหนึ่งวัน เก็บการ์ดนี้ไว้ ตอนเขียนโค้ดค่อยเปิดเทียบดู อีกสองสัปดาห์ก็คุ้นมือเอง:

\begin{center}
\begin{tabularx}{\linewidth}{|X|X|}
\hline
\textbf{ต้องการทำอะไร} & \textbf{วิธีเขียน} \tabularnewline
\hline
ประกาศตัวแปรแบบอ่านอย่างเดียว & \texttt{val x = 1} \tabularnewline
ประกาศตัวแปรแบบแก้ได้ & \texttt{var x = 1} \tabularnewline
ประกาศฟังก์ชัน & fun ชื่อ(พารามิเตอร์): ชนิดค่าที่ส่งกลับ \{ ... \} \tabularnewline
เช็กเงื่อนไข & if (เงื่อนไข) \{...\} else \{...\} \tabularnewline
หลายทางเลือก & when \{ เงื่อนไข -> ผลลัพธ์ \} \tabularnewline
วนลูป & for (x in ลิสต์) \{ ... \} \tabularnewline
ต่อข้อความ & "สวัสดี \$name" \tabularnewline
จัดการค่าที่อาจเป็น null & \texttt{x?.\allowbreak length ?: 0} \tabularnewline
ประกาศ data class & \texttt{data class X(val a: Int, val b: String)} \tabularnewline
\hline
\end{tabularx}
\end{center}

\begin{center}\small\color{gray}การ์ดสรุปไวยากรณ์ Kotlin (ฉบับมือใหม่)\end{center}

\begin{tcbpractice}{ลงมือทำด้วยกัน: เล่น Kotlin 10 นาที โดยไม่ต้องติดตั้งอะไรเลย}
เปิดเบราว์เซอร์ไปที่ play.kotlinlang.org (สนามซ้อมออนไลน์อย่างเป็นทางการของ Kotlin) แล้วเอาโค้ดแต่ละช่วงในบทนี้ไปวางรันดูทีละชุด ลองเปลี่ยนตัวเลขกับข้อความดูว่าผลลัพธ์เปลี่ยนยังไง 10 นาทีแบบ “เล่นก่อนเรียน” นี้ จะทำให้คุณเริ่มบทที่ 51 ได้ง่ายขึ้นเท่าตัว

\end{tcbpractice}

\begin{tcbsummary}{สรุปบทนี้}
\begin{itemize}
  \item val มาก่อน ใช้ var อย่างระวัง; ชนิดข้อมูลอนุมานให้เองได้ หรือประกาศชัดหลังทวิภาคก็ได้
  \item ฟังก์ชัน = ตรรกะที่ห่อไว้พร้อมใช้; สตริงเทมเพลต \$ คือของเด่นสะดวกสบาย
  \item if/when ไว้ตัดสินใจ for ไว้วนลูป; when คือ "ผู้รอบด้าน" ของงานพัฒนา Android
  \item null safety คือเครื่องรางกันตายของ Kotlin: ชนิดปกติปฏิเสธ null ส่วนชนิดที่เป็น null ได้ต้องจัดการด้วย ? / ?:
  \item data class คือกล่องใส่ข้อมูล หนึ่งฟิลด์ต่อหนึ่งบรรทัด โปรเจกต์จริงใช้ทุกวัน
  \item ไวยากรณ์ 20\% ทำงานได้ 80\%: ใช้เป็นก่อน แล้วค่อยไปลุ้นความชำนาญ
\end{itemize}

\end{tcbsummary}
